% economic_model.tex — Section 2: Economic Model of Crop Waste Contagion
% Included via % economic_model.tex — Section 2: Economic Model of Crop Waste Contagion
% Included via % economic_model.tex — Section 2: Economic Model of Crop Waste Contagion
% Included via % economic_model.tex — Section 2: Economic Model of Crop Waste Contagion
% Included via \input{sections/economic_model} in main.tex
% The \section and \label are in main.tex, not here.

Why does a bumper harvest in one county cause crop waste hundreds of miles away?
Standard agronomic models treat waste as a local phenomenon---drought, pest damage, frost---and ignore the market channel entirely.
We formalize the economic mechanism that propagates waste across regions, derive testable predictions, and show how these predictions map to specific architectural choices in our graph neural network.

\subsection{Model Setup}
\label{sec:model_setup}

Consider $N$ agricultural regions (counties) indexed by $i = 1, \ldots, N$.
Each region $i$ grows a commodity with production quantity $q_i$, determined by biophysical conditions (weather, soil, planted acreage), and faces a harvesting cost $c_i > 0$ encompassing labor, fuel, and equipment.
The farmer in region $i$ observes a local market price $p_i$ and harvests if and only if revenue exceeds cost:
\begin{equation}
\label{eq:harvest_condition}
p_i \cdot q_i > c_i.
\end{equation}
When this condition fails, the standing crop is abandoned---left unharvested and wasted.

The local price depends on aggregate regional supply through a linear inverse demand system:
\begin{equation}
\label{eq:price}
p_i = P^{\mathrm{global}} - \sum_{j=1}^{N} \beta_{ij} \, q_j,
\end{equation}
where $P^{\mathrm{global}} > 0$ is the national or world reference price and $\beta_{ij} \geq 0$ captures \emph{market connectivity} between regions $i$ and $j$.

\paragraph{Interpretation of $\beta_{ij}$.}
The coefficient $\beta_{ij}$ encodes how strongly production in region $j$ depresses the price received by farmers in region $i$.
Four properties govern its structure:
\begin{enumerate}[label=(\roman*)]
    \item $\beta_{ij} > 0$ if and only if regions $i$ and $j$ compete in the same downstream market (shared processors, packers, or distribution channels).
    \item $\beta_{ij}$ increases with geographic proximity, since transportation costs determine which buyers are accessible.
    \item $\beta_{ij}$ increases with commodity similarity---two adjacent corn-producing counties interact more strongly than a corn county and a soybean county.
    \item $\beta_{ij} \approx 0$ for distant or market-disconnected regions.
\end{enumerate}
These properties make the matrix $\mathbf{B} = [\beta_{ij}]$ sparse, spatially structured, and commodity-dependent---precisely the adjacency structure our graph model learns from data.

\paragraph{Waste probability.}
Define the \emph{profit margin} for region $i$ as
\begin{equation}
\label{eq:margin}
\pi_i \;=\; p_i \cdot q_i - c_i \;=\; \left(P^{\mathrm{global}} - \sum_{j} \beta_{ij}\, q_j \right) q_i \;-\; c_i.
\end{equation}
Region $i$ wastes its crop whenever $\pi_i \leq 0$.
In a stochastic setting where $q_j$ is drawn from a distribution $F_j$ (reflecting weather uncertainty), the waste probability is
\begin{equation}
\label{eq:waste_prob}
\Pr(\text{waste}_i) = \Pr\!\left(\pi_i \leq 0\right) = \Pr\!\left(\sum_{j} \beta_{ij}\, q_j \geq P^{\mathrm{global}} - \frac{c_i}{q_i}\right).
\end{equation}

\subsection{Waste Contagion}

\begin{proposition}[Waste Contagion]
\label{prop:contagion}
Let region $j$ experience a positive supply shock: $q_j$ increases to $q_j + \Delta q_j$ with $\Delta q_j > 0$, while all other quantities $\{q_k\}_{k \neq j}$ remain unchanged.
If $\beta_{ij} > 0$, then the waste probability in region $i$ strictly increases, even when region $i$'s own biophysical conditions are unchanged.
\end{proposition}

\begin{proof}
From Equation~\eqref{eq:price}, the price in region $i$ after the shock is
\[
p_i' = P^{\mathrm{global}} - \sum_{k \neq j} \beta_{ik}\, q_k - \beta_{ij}(q_j + \Delta q_j) = p_i - \beta_{ij}\, \Delta q_j.
\]
Since $\beta_{ij} > 0$ and $\Delta q_j > 0$, we have $p_i' < p_i$.
The post-shock profit margin becomes
\[
\pi_i' = p_i' \cdot q_i - c_i = \pi_i - \beta_{ij}\, \Delta q_j \cdot q_i < \pi_i.
\]
Because $q_i > 0$ (region $i$ has a standing crop), the margin drops by the strictly positive quantity $\beta_{ij}\, \Delta q_j \, q_i$.
Hence $\Pr(\pi_i' \leq 0) > \Pr(\pi_i \leq 0)$, and the waste probability in region $i$ strictly increases.
\end{proof}

Proposition~\ref{prop:contagion} captures the central paradox of crop waste: good weather in one county can destroy value in another.
The mechanism is purely economic---no pathogen spreads, no pest migrates.
A bumper harvest floods shared markets, crashes the local price, and pushes marginal farmers below their break-even threshold.
The strength of this contagion is governed entirely by $\beta_{ij}$, the market connectivity~\citep{acemoglu2015systemic}.

\begin{corollary}[Asymmetric Propagation]
\label{cor:asymmetry}
Positive and negative supply shocks in region $j$ have asymmetric effects on waste in region $i$:
\begin{enumerate}[label=(\alph*)]
    \item \textbf{Bumper harvest} ($\Delta q_j > 0$): $p_i$ falls by $\beta_{ij}\, \Delta q_j$, \emph{increasing} waste probability in $i$.
    \item \textbf{Crop failure} ($\Delta q_j < 0$): $p_i$ rises by $\beta_{ij}\, |\Delta q_j|$, \emph{decreasing} waste probability in $i$.
\end{enumerate}
Consequently, the effect of neighbor shocks on waste risk is directionally asymmetric: positive shocks are harmful and negative shocks are beneficial to region $i$'s harvest viability.
\end{corollary}

The proof follows immediately from Proposition~\ref{prop:contagion} by substituting $\Delta q_j < 0$.
This asymmetry has a direct architectural implication: a predictive model must condition on the \emph{sign} of neighbor shocks, not merely their magnitude.
Symmetric message-passing---treating $+\Delta q_j$ and $-\Delta q_j$ identically---discards the directional information that distinguishes waste-inducing from waste-relieving shocks.
This motivates the directional shock conditioning in our Economic Cascade Message Passing (ECMP) layer.

\subsection{Testable Hypotheses}
\label{sec:hypotheses}

The model generates three sharp predictions that map directly to ablation experiments.
Each hypothesis corresponds to a specific comparison row in our results table.

\begin{hypothesis}[Graph Structure Improves Prediction]
\label{hyp:graph}
Because waste contagion propagates through the connectivity matrix $\mathbf{B}$, models that capture inter-regional dependencies (graph-based) will outperform models that treat each region independently.
\end{hypothesis}

\noindent\textit{Test.}\; Compare the full ECMP model (Row~5) against region-independent baselines---gradient-boosted trees and standalone neural networks (Rows~1--2).
If $\mathbf{B}$ is informative, graph models should achieve lower error on out-of-sample waste prediction.

\begin{hypothesis}[Economic Edges Outperform Geographic Proximity]
\label{hyp:edges}
Since $\beta_{ij}$ depends on market structure, commodity similarity, and transportation networks---not geographic distance alone---edges derived from economic connectivity will carry more predictive signal than edges based purely on spatial adjacency.
\end{hypothesis}

\noindent\textit{Test.}\; Compare ECMP with learned economic edges (Row~5) against a geographic-only graph model using $k$-nearest spatial neighbors (Row~3).
A significant gap confirms that market structure, not mere proximity, drives contagion.

\begin{hypothesis}[Directional Conditioning Captures Asymmetry]
\label{hyp:asymmetry}
By Corollary~\ref{cor:asymmetry}, positive and negative neighbor shocks have opposite effects on waste risk.
Models that condition on shock direction---distinguishing surplus from deficit signals---will outperform symmetric alternatives that aggregate neighbor information without regard to sign.
\end{hypothesis}

\noindent\textit{Test.}\; Compare ECMP with directional shock conditioning (Row~5) against a symmetric message-passing variant (Row~4).
If the asymmetry from Corollary~\ref{cor:asymmetry} matters empirically, directional conditioning should yield measurable gains.

\paragraph{From theory to architecture.}
These three hypotheses jointly motivate every key design choice in ECMP\@.
Hypothesis~\ref{hyp:graph} justifies graph neural networks over tabular baselines.
Hypothesis~\ref{hyp:edges} drives our economic edge construction---using commodity flows, processor locations, and trade data rather than Euclidean distance.
Hypothesis~\ref{hyp:asymmetry} motivates the signed message-passing mechanism that separates surplus and deficit channels during neighborhood aggregation.
Section~\ref{sec:methods} details how we translate each theoretical requirement into a concrete architectural component.
 in main.tex
% The \section and \label are in main.tex, not here.

Why does a bumper harvest in one county cause crop waste hundreds of miles away?
Standard agronomic models treat waste as a local phenomenon---drought, pest damage, frost---and ignore the market channel entirely.
We formalize the economic mechanism that propagates waste across regions, derive testable predictions, and show how these predictions map to specific architectural choices in our graph neural network.

\subsection{Model Setup}
\label{sec:model_setup}

Consider $N$ agricultural regions (counties) indexed by $i = 1, \ldots, N$.
Each region $i$ grows a commodity with production quantity $q_i$, determined by biophysical conditions (weather, soil, planted acreage), and faces a harvesting cost $c_i > 0$ encompassing labor, fuel, and equipment.
The farmer in region $i$ observes a local market price $p_i$ and harvests if and only if revenue exceeds cost:
\begin{equation}
\label{eq:harvest_condition}
p_i \cdot q_i > c_i.
\end{equation}
When this condition fails, the standing crop is abandoned---left unharvested and wasted.

The local price depends on aggregate regional supply through a linear inverse demand system:
\begin{equation}
\label{eq:price}
p_i = P^{\mathrm{global}} - \sum_{j=1}^{N} \beta_{ij} \, q_j,
\end{equation}
where $P^{\mathrm{global}} > 0$ is the national or world reference price and $\beta_{ij} \geq 0$ captures \emph{market connectivity} between regions $i$ and $j$.

\paragraph{Interpretation of $\beta_{ij}$.}
The coefficient $\beta_{ij}$ encodes how strongly production in region $j$ depresses the price received by farmers in region $i$.
Four properties govern its structure:
\begin{enumerate}[label=(\roman*)]
    \item $\beta_{ij} > 0$ if and only if regions $i$ and $j$ compete in the same downstream market (shared processors, packers, or distribution channels).
    \item $\beta_{ij}$ increases with geographic proximity, since transportation costs determine which buyers are accessible.
    \item $\beta_{ij}$ increases with commodity similarity---two adjacent corn-producing counties interact more strongly than a corn county and a soybean county.
    \item $\beta_{ij} \approx 0$ for distant or market-disconnected regions.
\end{enumerate}
These properties make the matrix $\mathbf{B} = [\beta_{ij}]$ sparse, spatially structured, and commodity-dependent---precisely the adjacency structure our graph model learns from data.

\paragraph{Waste probability.}
Define the \emph{profit margin} for region $i$ as
\begin{equation}
\label{eq:margin}
\pi_i \;=\; p_i \cdot q_i - c_i \;=\; \left(P^{\mathrm{global}} - \sum_{j} \beta_{ij}\, q_j \right) q_i \;-\; c_i.
\end{equation}
Region $i$ wastes its crop whenever $\pi_i \leq 0$.
In a stochastic setting where $q_j$ is drawn from a distribution $F_j$ (reflecting weather uncertainty), the waste probability is
\begin{equation}
\label{eq:waste_prob}
\Pr(\text{waste}_i) = \Pr\!\left(\pi_i \leq 0\right) = \Pr\!\left(\sum_{j} \beta_{ij}\, q_j \geq P^{\mathrm{global}} - \frac{c_i}{q_i}\right).
\end{equation}

\subsection{Waste Contagion}

\begin{proposition}[Waste Contagion]
\label{prop:contagion}
Let region $j$ experience a positive supply shock: $q_j$ increases to $q_j + \Delta q_j$ with $\Delta q_j > 0$, while all other quantities $\{q_k\}_{k \neq j}$ remain unchanged.
If $\beta_{ij} > 0$, then the waste probability in region $i$ strictly increases, even when region $i$'s own biophysical conditions are unchanged.
\end{proposition}

\begin{proof}
From Equation~\eqref{eq:price}, the price in region $i$ after the shock is
\[
p_i' = P^{\mathrm{global}} - \sum_{k \neq j} \beta_{ik}\, q_k - \beta_{ij}(q_j + \Delta q_j) = p_i - \beta_{ij}\, \Delta q_j.
\]
Since $\beta_{ij} > 0$ and $\Delta q_j > 0$, we have $p_i' < p_i$.
The post-shock profit margin becomes
\[
\pi_i' = p_i' \cdot q_i - c_i = \pi_i - \beta_{ij}\, \Delta q_j \cdot q_i < \pi_i.
\]
Because $q_i > 0$ (region $i$ has a standing crop), the margin drops by the strictly positive quantity $\beta_{ij}\, \Delta q_j \, q_i$.
Hence $\Pr(\pi_i' \leq 0) > \Pr(\pi_i \leq 0)$, and the waste probability in region $i$ strictly increases.
\end{proof}

Proposition~\ref{prop:contagion} captures the central paradox of crop waste: good weather in one county can destroy value in another.
The mechanism is purely economic---no pathogen spreads, no pest migrates.
A bumper harvest floods shared markets, crashes the local price, and pushes marginal farmers below their break-even threshold.
The strength of this contagion is governed entirely by $\beta_{ij}$, the market connectivity~\citep{acemoglu2015systemic}.

\begin{corollary}[Asymmetric Propagation]
\label{cor:asymmetry}
Positive and negative supply shocks in region $j$ have asymmetric effects on waste in region $i$:
\begin{enumerate}[label=(\alph*)]
    \item \textbf{Bumper harvest} ($\Delta q_j > 0$): $p_i$ falls by $\beta_{ij}\, \Delta q_j$, \emph{increasing} waste probability in $i$.
    \item \textbf{Crop failure} ($\Delta q_j < 0$): $p_i$ rises by $\beta_{ij}\, |\Delta q_j|$, \emph{decreasing} waste probability in $i$.
\end{enumerate}
Consequently, the effect of neighbor shocks on waste risk is directionally asymmetric: positive shocks are harmful and negative shocks are beneficial to region $i$'s harvest viability.
\end{corollary}

The proof follows immediately from Proposition~\ref{prop:contagion} by substituting $\Delta q_j < 0$.
This asymmetry has a direct architectural implication: a predictive model must condition on the \emph{sign} of neighbor shocks, not merely their magnitude.
Symmetric message-passing---treating $+\Delta q_j$ and $-\Delta q_j$ identically---discards the directional information that distinguishes waste-inducing from waste-relieving shocks.
This motivates the directional shock conditioning in our Economic Cascade Message Passing (ECMP) layer.

\subsection{Testable Hypotheses}
\label{sec:hypotheses}

The model generates three sharp predictions that map directly to ablation experiments.
Each hypothesis corresponds to a specific comparison row in our results table.

\begin{hypothesis}[Graph Structure Improves Prediction]
\label{hyp:graph}
Because waste contagion propagates through the connectivity matrix $\mathbf{B}$, models that capture inter-regional dependencies (graph-based) will outperform models that treat each region independently.
\end{hypothesis}

\noindent\textit{Test.}\; Compare the full ECMP model (Row~5) against region-independent baselines---gradient-boosted trees and standalone neural networks (Rows~1--2).
If $\mathbf{B}$ is informative, graph models should achieve lower error on out-of-sample waste prediction.

\begin{hypothesis}[Economic Edges Outperform Geographic Proximity]
\label{hyp:edges}
Since $\beta_{ij}$ depends on market structure, commodity similarity, and transportation networks---not geographic distance alone---edges derived from economic connectivity will carry more predictive signal than edges based purely on spatial adjacency.
\end{hypothesis}

\noindent\textit{Test.}\; Compare ECMP with learned economic edges (Row~5) against a geographic-only graph model using $k$-nearest spatial neighbors (Row~3).
A significant gap confirms that market structure, not mere proximity, drives contagion.

\begin{hypothesis}[Directional Conditioning Captures Asymmetry]
\label{hyp:asymmetry}
By Corollary~\ref{cor:asymmetry}, positive and negative neighbor shocks have opposite effects on waste risk.
Models that condition on shock direction---distinguishing surplus from deficit signals---will outperform symmetric alternatives that aggregate neighbor information without regard to sign.
\end{hypothesis}

\noindent\textit{Test.}\; Compare ECMP with directional shock conditioning (Row~5) against a symmetric message-passing variant (Row~4).
If the asymmetry from Corollary~\ref{cor:asymmetry} matters empirically, directional conditioning should yield measurable gains.

\paragraph{From theory to architecture.}
These three hypotheses jointly motivate every key design choice in ECMP\@.
Hypothesis~\ref{hyp:graph} justifies graph neural networks over tabular baselines.
Hypothesis~\ref{hyp:edges} drives our economic edge construction---using commodity flows, processor locations, and trade data rather than Euclidean distance.
Hypothesis~\ref{hyp:asymmetry} motivates the signed message-passing mechanism that separates surplus and deficit channels during neighborhood aggregation.
Section~\ref{sec:methods} details how we translate each theoretical requirement into a concrete architectural component.
 in main.tex
% The \section and \label are in main.tex, not here.

Why does a bumper harvest in one county cause crop waste hundreds of miles away?
Standard agronomic models treat waste as a local phenomenon---drought, pest damage, frost---and ignore the market channel entirely.
We formalize the economic mechanism that propagates waste across regions, derive testable predictions, and show how these predictions map to specific architectural choices in our graph neural network.

\subsection{Model Setup}
\label{sec:model_setup}

Consider $N$ agricultural regions (counties) indexed by $i = 1, \ldots, N$.
Each region $i$ grows a commodity with production quantity $q_i$, determined by biophysical conditions (weather, soil, planted acreage), and faces a harvesting cost $c_i > 0$ encompassing labor, fuel, and equipment.
The farmer in region $i$ observes a local market price $p_i$ and harvests if and only if revenue exceeds cost:
\begin{equation}
\label{eq:harvest_condition}
p_i \cdot q_i > c_i.
\end{equation}
When this condition fails, the standing crop is abandoned---left unharvested and wasted.

The local price depends on aggregate regional supply through a linear inverse demand system:
\begin{equation}
\label{eq:price}
p_i = P^{\mathrm{global}} - \sum_{j=1}^{N} \beta_{ij} \, q_j,
\end{equation}
where $P^{\mathrm{global}} > 0$ is the national or world reference price and $\beta_{ij} \geq 0$ captures \emph{market connectivity} between regions $i$ and $j$.

\paragraph{Interpretation of $\beta_{ij}$.}
The coefficient $\beta_{ij}$ encodes how strongly production in region $j$ depresses the price received by farmers in region $i$.
Four properties govern its structure:
\begin{enumerate}[label=(\roman*)]
    \item $\beta_{ij} > 0$ if and only if regions $i$ and $j$ compete in the same downstream market (shared processors, packers, or distribution channels).
    \item $\beta_{ij}$ increases with geographic proximity, since transportation costs determine which buyers are accessible.
    \item $\beta_{ij}$ increases with commodity similarity---two adjacent corn-producing counties interact more strongly than a corn county and a soybean county.
    \item $\beta_{ij} \approx 0$ for distant or market-disconnected regions.
\end{enumerate}
These properties make the matrix $\mathbf{B} = [\beta_{ij}]$ sparse, spatially structured, and commodity-dependent---precisely the adjacency structure our graph model learns from data.

\paragraph{Waste probability.}
Define the \emph{profit margin} for region $i$ as
\begin{equation}
\label{eq:margin}
\pi_i \;=\; p_i \cdot q_i - c_i \;=\; \left(P^{\mathrm{global}} - \sum_{j} \beta_{ij}\, q_j \right) q_i \;-\; c_i.
\end{equation}
Region $i$ wastes its crop whenever $\pi_i \leq 0$.
In a stochastic setting where $q_j$ is drawn from a distribution $F_j$ (reflecting weather uncertainty), the waste probability is
\begin{equation}
\label{eq:waste_prob}
\Pr(\text{waste}_i) = \Pr\!\left(\pi_i \leq 0\right) = \Pr\!\left(\sum_{j} \beta_{ij}\, q_j \geq P^{\mathrm{global}} - \frac{c_i}{q_i}\right).
\end{equation}

\subsection{Waste Contagion}

\begin{proposition}[Waste Contagion]
\label{prop:contagion}
Let region $j$ experience a positive supply shock: $q_j$ increases to $q_j + \Delta q_j$ with $\Delta q_j > 0$, while all other quantities $\{q_k\}_{k \neq j}$ remain unchanged.
If $\beta_{ij} > 0$, then the waste probability in region $i$ strictly increases, even when region $i$'s own biophysical conditions are unchanged.
\end{proposition}

\begin{proof}
From Equation~\eqref{eq:price}, the price in region $i$ after the shock is
\[
p_i' = P^{\mathrm{global}} - \sum_{k \neq j} \beta_{ik}\, q_k - \beta_{ij}(q_j + \Delta q_j) = p_i - \beta_{ij}\, \Delta q_j.
\]
Since $\beta_{ij} > 0$ and $\Delta q_j > 0$, we have $p_i' < p_i$.
The post-shock profit margin becomes
\[
\pi_i' = p_i' \cdot q_i - c_i = \pi_i - \beta_{ij}\, \Delta q_j \cdot q_i < \pi_i.
\]
Because $q_i > 0$ (region $i$ has a standing crop), the margin drops by the strictly positive quantity $\beta_{ij}\, \Delta q_j \, q_i$.
Hence $\Pr(\pi_i' \leq 0) > \Pr(\pi_i \leq 0)$, and the waste probability in region $i$ strictly increases.
\end{proof}

Proposition~\ref{prop:contagion} captures the central paradox of crop waste: good weather in one county can destroy value in another.
The mechanism is purely economic---no pathogen spreads, no pest migrates.
A bumper harvest floods shared markets, crashes the local price, and pushes marginal farmers below their break-even threshold.
The strength of this contagion is governed entirely by $\beta_{ij}$, the market connectivity~\citep{acemoglu2015systemic}.

\begin{corollary}[Asymmetric Propagation]
\label{cor:asymmetry}
Positive and negative supply shocks in region $j$ have asymmetric effects on waste in region $i$:
\begin{enumerate}[label=(\alph*)]
    \item \textbf{Bumper harvest} ($\Delta q_j > 0$): $p_i$ falls by $\beta_{ij}\, \Delta q_j$, \emph{increasing} waste probability in $i$.
    \item \textbf{Crop failure} ($\Delta q_j < 0$): $p_i$ rises by $\beta_{ij}\, |\Delta q_j|$, \emph{decreasing} waste probability in $i$.
\end{enumerate}
Consequently, the effect of neighbor shocks on waste risk is directionally asymmetric: positive shocks are harmful and negative shocks are beneficial to region $i$'s harvest viability.
\end{corollary}

The proof follows immediately from Proposition~\ref{prop:contagion} by substituting $\Delta q_j < 0$.
This asymmetry has a direct architectural implication: a predictive model must condition on the \emph{sign} of neighbor shocks, not merely their magnitude.
Symmetric message-passing---treating $+\Delta q_j$ and $-\Delta q_j$ identically---discards the directional information that distinguishes waste-inducing from waste-relieving shocks.
This motivates the directional shock conditioning in our Economic Cascade Message Passing (ECMP) layer.

\subsection{Testable Hypotheses}
\label{sec:hypotheses}

The model generates three sharp predictions that map directly to ablation experiments.
Each hypothesis corresponds to a specific comparison row in our results table.

\begin{hypothesis}[Graph Structure Improves Prediction]
\label{hyp:graph}
Because waste contagion propagates through the connectivity matrix $\mathbf{B}$, models that capture inter-regional dependencies (graph-based) will outperform models that treat each region independently.
\end{hypothesis}

\noindent\textit{Test.}\; Compare the full ECMP model (Row~5) against region-independent baselines---gradient-boosted trees and standalone neural networks (Rows~1--2).
If $\mathbf{B}$ is informative, graph models should achieve lower error on out-of-sample waste prediction.

\begin{hypothesis}[Economic Edges Outperform Geographic Proximity]
\label{hyp:edges}
Since $\beta_{ij}$ depends on market structure, commodity similarity, and transportation networks---not geographic distance alone---edges derived from economic connectivity will carry more predictive signal than edges based purely on spatial adjacency.
\end{hypothesis}

\noindent\textit{Test.}\; Compare ECMP with learned economic edges (Row~5) against a geographic-only graph model using $k$-nearest spatial neighbors (Row~3).
A significant gap confirms that market structure, not mere proximity, drives contagion.

\begin{hypothesis}[Directional Conditioning Captures Asymmetry]
\label{hyp:asymmetry}
By Corollary~\ref{cor:asymmetry}, positive and negative neighbor shocks have opposite effects on waste risk.
Models that condition on shock direction---distinguishing surplus from deficit signals---will outperform symmetric alternatives that aggregate neighbor information without regard to sign.
\end{hypothesis}

\noindent\textit{Test.}\; Compare ECMP with directional shock conditioning (Row~5) against a symmetric message-passing variant (Row~4).
If the asymmetry from Corollary~\ref{cor:asymmetry} matters empirically, directional conditioning should yield measurable gains.

\paragraph{From theory to architecture.}
These three hypotheses jointly motivate every key design choice in ECMP\@.
Hypothesis~\ref{hyp:graph} justifies graph neural networks over tabular baselines.
Hypothesis~\ref{hyp:edges} drives our economic edge construction---using commodity flows, processor locations, and trade data rather than Euclidean distance.
Hypothesis~\ref{hyp:asymmetry} motivates the signed message-passing mechanism that separates surplus and deficit channels during neighborhood aggregation.
Section~\ref{sec:methods} details how we translate each theoretical requirement into a concrete architectural component.
 in main.tex
% The \section and \label are in main.tex, not here.

Why does a bumper harvest in one county cause crop waste hundreds of miles away?
Standard agronomic models treat waste as a local phenomenon---drought, pest damage, frost---and ignore the market channel entirely.
We formalize the economic mechanism that propagates waste across regions, derive testable predictions, and show how these predictions map to specific architectural choices in our graph neural network.

\subsection{Model Setup}
\label{sec:model_setup}

Consider $N$ agricultural regions (counties) indexed by $i = 1, \ldots, N$.
Each region $i$ grows a commodity with production quantity $q_i$, determined by biophysical conditions (weather, soil, planted acreage), and faces a harvesting cost $c_i > 0$ encompassing labor, fuel, and equipment.
The farmer in region $i$ observes a local market price $p_i$ and harvests if and only if revenue exceeds cost:
\begin{equation}
\label{eq:harvest_condition}
p_i \cdot q_i > c_i.
\end{equation}
When this condition fails, the standing crop is abandoned---left unharvested and wasted.

The local price depends on aggregate regional supply through a linear inverse demand system:
\begin{equation}
\label{eq:price}
p_i = P^{\mathrm{global}} - \sum_{j=1}^{N} \beta_{ij} \, q_j,
\end{equation}
where $P^{\mathrm{global}} > 0$ is the national or world reference price and $\beta_{ij} \geq 0$ captures \emph{market connectivity} between regions $i$ and $j$.

\paragraph{Interpretation of $\beta_{ij}$.}
The coefficient $\beta_{ij}$ encodes how strongly production in region $j$ depresses the price received by farmers in region $i$.
Four properties govern its structure:
\begin{enumerate}[label=(\roman*)]
    \item $\beta_{ij} > 0$ if and only if regions $i$ and $j$ compete in the same downstream market (shared processors, packers, or distribution channels).
    \item $\beta_{ij}$ increases with geographic proximity, since transportation costs determine which buyers are accessible.
    \item $\beta_{ij}$ increases with commodity similarity---two adjacent corn-producing counties interact more strongly than a corn county and a soybean county.
    \item $\beta_{ij} \approx 0$ for distant or market-disconnected regions.
\end{enumerate}
These properties make the matrix $\mathbf{B} = [\beta_{ij}]$ sparse, spatially structured, and commodity-dependent---precisely the adjacency structure our graph model learns from data.

\paragraph{Waste probability.}
Define the \emph{profit margin} for region $i$ as
\begin{equation}
\label{eq:margin}
\pi_i \;=\; p_i \cdot q_i - c_i \;=\; \left(P^{\mathrm{global}} - \sum_{j} \beta_{ij}\, q_j \right) q_i \;-\; c_i.
\end{equation}
Region $i$ wastes its crop whenever $\pi_i \leq 0$.
In a stochastic setting where $q_j$ is drawn from a distribution $F_j$ (reflecting weather uncertainty), the waste probability is
\begin{equation}
\label{eq:waste_prob}
\Pr(\text{waste}_i) = \Pr\!\left(\pi_i \leq 0\right) = \Pr\!\left(\sum_{j} \beta_{ij}\, q_j \geq P^{\mathrm{global}} - \frac{c_i}{q_i}\right).
\end{equation}

\subsection{Waste Contagion}

\begin{proposition}[Waste Contagion]
\label{prop:contagion}
Let region $j$ experience a positive supply shock: $q_j$ increases to $q_j + \Delta q_j$ with $\Delta q_j > 0$, while all other quantities $\{q_k\}_{k \neq j}$ remain unchanged.
If $\beta_{ij} > 0$, then the waste probability in region $i$ strictly increases, even when region $i$'s own biophysical conditions are unchanged.
\end{proposition}

\begin{proof}
From Equation~\eqref{eq:price}, the price in region $i$ after the shock is
\[
p_i' = P^{\mathrm{global}} - \sum_{k \neq j} \beta_{ik}\, q_k - \beta_{ij}(q_j + \Delta q_j) = p_i - \beta_{ij}\, \Delta q_j.
\]
Since $\beta_{ij} > 0$ and $\Delta q_j > 0$, we have $p_i' < p_i$.
The post-shock profit margin becomes
\[
\pi_i' = p_i' \cdot q_i - c_i = \pi_i - \beta_{ij}\, \Delta q_j \cdot q_i < \pi_i.
\]
Because $q_i > 0$ (region $i$ has a standing crop), the margin drops by the strictly positive quantity $\beta_{ij}\, \Delta q_j \, q_i$.
Hence $\Pr(\pi_i' \leq 0) > \Pr(\pi_i \leq 0)$, and the waste probability in region $i$ strictly increases.
\end{proof}

Proposition~\ref{prop:contagion} captures the central paradox of crop waste: good weather in one county can destroy value in another.
The mechanism is purely economic---no pathogen spreads, no pest migrates.
A bumper harvest floods shared markets, crashes the local price, and pushes marginal farmers below their break-even threshold.
The strength of this contagion is governed entirely by $\beta_{ij}$, the market connectivity~\citep{acemoglu2015systemic}.

\begin{corollary}[Asymmetric Propagation]
\label{cor:asymmetry}
Positive and negative supply shocks in region $j$ have asymmetric effects on waste in region $i$:
\begin{enumerate}[label=(\alph*)]
    \item \textbf{Bumper harvest} ($\Delta q_j > 0$): $p_i$ falls by $\beta_{ij}\, \Delta q_j$, \emph{increasing} waste probability in $i$.
    \item \textbf{Crop failure} ($\Delta q_j < 0$): $p_i$ rises by $\beta_{ij}\, |\Delta q_j|$, \emph{decreasing} waste probability in $i$.
\end{enumerate}
Consequently, the effect of neighbor shocks on waste risk is directionally asymmetric: positive shocks are harmful and negative shocks are beneficial to region $i$'s harvest viability.
\end{corollary}

The proof follows immediately from Proposition~\ref{prop:contagion} by substituting $\Delta q_j < 0$.
This asymmetry has a direct architectural implication: a predictive model must condition on the \emph{sign} of neighbor shocks, not merely their magnitude.
Symmetric message-passing---treating $+\Delta q_j$ and $-\Delta q_j$ identically---discards the directional information that distinguishes waste-inducing from waste-relieving shocks.
This motivates the directional shock conditioning in our Economic Cascade Message Passing (ECMP) layer.

\subsection{Testable Hypotheses}
\label{sec:hypotheses}

The model generates three sharp predictions that map directly to ablation experiments.
Each hypothesis corresponds to a specific comparison row in our results table.

\begin{hypothesis}[Graph Structure Improves Prediction]
\label{hyp:graph}
Because waste contagion propagates through the connectivity matrix $\mathbf{B}$, models that capture inter-regional dependencies (graph-based) will outperform models that treat each region independently.
\end{hypothesis}

\noindent\textit{Test.}\; Compare the full ECMP model (Row~5) against region-independent baselines---gradient-boosted trees and standalone neural networks (Rows~1--2).
If $\mathbf{B}$ is informative, graph models should achieve lower error on out-of-sample waste prediction.

\begin{hypothesis}[Economic Edges Outperform Geographic Proximity]
\label{hyp:edges}
Since $\beta_{ij}$ depends on market structure, commodity similarity, and transportation networks---not geographic distance alone---edges derived from economic connectivity will carry more predictive signal than edges based purely on spatial adjacency.
\end{hypothesis}

\noindent\textit{Test.}\; Compare ECMP with learned economic edges (Row~5) against a geographic-only graph model using $k$-nearest spatial neighbors (Row~3).
A significant gap confirms that market structure, not mere proximity, drives contagion.

\begin{hypothesis}[Directional Conditioning Captures Asymmetry]
\label{hyp:asymmetry}
By Corollary~\ref{cor:asymmetry}, positive and negative neighbor shocks have opposite effects on waste risk.
Models that condition on shock direction---distinguishing surplus from deficit signals---will outperform symmetric alternatives that aggregate neighbor information without regard to sign.
\end{hypothesis}

\noindent\textit{Test.}\; Compare ECMP with directional shock conditioning (Row~5) against a symmetric message-passing variant (Row~4).
If the asymmetry from Corollary~\ref{cor:asymmetry} matters empirically, directional conditioning should yield measurable gains.

\paragraph{From theory to architecture.}
These three hypotheses jointly motivate every key design choice in ECMP\@.
Hypothesis~\ref{hyp:graph} justifies graph neural networks over tabular baselines.
Hypothesis~\ref{hyp:edges} drives our economic edge construction---using commodity flows, processor locations, and trade data rather than Euclidean distance.
Hypothesis~\ref{hyp:asymmetry} motivates the signed message-passing mechanism that separates surplus and deficit channels during neighborhood aggregation.
Section~\ref{sec:methods} details how we translate each theoretical requirement into a concrete architectural component.
